\section*{Bab \ref{ch:b3:microbiomes} --- Mikrobiom dan Simbiosis}

\begin{solution}{exo:b3:microbiomes:1}
\omterm{def:b3:microbiomes:symbiosis}{Mutualisme}: keduanya beroleh untung --- kacang dan rizobium, karang dan alga.
\omterm{def:b3:microbiomes:symbiosis}{Komensalisme}: satu beroleh untung, yang lain tidak terpengaruh --- kebanyakan
bakteri usus yang makan dari apa yang tak dapat dicerna inangnya. Parasitisme:
satu beroleh untung dengan ongkos yang lain --- \emph{Wolbachia} yang membunuh
embrio jantan. Pasangan yang sama dapat bergeser: sebuah komensal usus yang
menyeberangi dindingnya menjadi patogen, sebuah rizobium yang berhenti
memfiksasi menjadi parasit bintilnya, dan sebuah alga di bawah tekanan panas
merugikan karang yang merumahkannya.
\end{solution}

\begin{solution}{exo:b3:microbiomes:2}
$H = -(0.5\ln 0.5 + 0.3\ln 0.3 + 0.1\ln 0.1 + 2\times 0.05\ln 0.05) =
0.347 + 0.361 + 0.230 + 0.300 = 1.24$; $e^{H} = 3.4$ spesies
mujarab. $D = 0.25 + 0.09 + 0.01 + 0.0025 + 0.0025 = 0.355$; $1/D =
2.8$.
\end{solution}

\begin{solution}{exo:b3:microbiomes:3}
Survei 16S melaporkan siapa yang ada di sana, sampai kira-kira marganya, dan
cacah \omterm{def:b3:genomics:ngs}{bacaan} nisbinya; ia tidak mengatakan apa pun tentang gen yang
diusungnya, melewatkan jamur dan \omterm{def:b3:virology:virus}{virus} (pemula lain), dan berat sebelah karena
cacah salinan dan pemula. Metagenomika senapan melaporkan gen, fungsi, dan,
dengan kedalaman yang cukup, genom serta galur utuh; ongkosnya lebih besar tiap
contoh, ia terbanjiri DNA inang pada contoh jaringan, dan ia tetap tak dapat
mengatakan gen mana yang terungkap atau sel mana yang hidup.
\end{solution}

\begin{solution}{exo:b3:microbiomes:4}
Meragikan serat menjadi \omterm{def:b3:microbiomes:gut}{asam lemak rantai pendek} yang memberi makan kolon dan
inangnya; mensintesis vitamin~K dan vitamin B; mengubah asam
empedu dan obat; \omterm{def:b3:microbiomes:gut}{ketahanan penjajahan} terhadap patogen; mendidik
sistem imun. Kekacauannya: kolitis \emph{Clostridioides difficile}
sesudah \omterm{def:b3:bacteriology:antibiotics}{antibiotik} (yang juga terlibat: penyakit radang usus,
kegemukan).
\end{solution}

\begin{solution}{exo:b3:microbiomes:5}
$\binom{20}{5} = \num{15504}$. Spesies dengan $10$: $1 - \binom{10}{5}/
\num{15504} = 1 - 252/\num{15504} = 0.984$; dengan $6$: $1 - 2002/\num{15504}
= 0.871$; dengan $3$: $1 - 6188/\num{15504} = 0.601$; dengan $1$: $1 -
\num{11628}/\num{15504} = 0.25$. $E[S_{5}] = 2.7$ spesies. Cakupannya: satu
tunggalan, $1 - 1/20 = 0.95$.
\end{solution}

\begin{solution}{exo:b3:microbiomes:6}
$400\times 10^{11} = 4\times 10^{13}$ bakteri, sekitar \qty{40}{g}. Gennya:
$1000$ spesies $\times$ \num{4000}, dikurangi tindihannya --- berorde
$3\times 10^{6}$ gen berbeda, lebih dari seratus kali lipat \num{20000}
milik manusia.
\end{solution}

\begin{solution}{exo:b3:microbiomes:7}
\omterm{def:b3:microbiomes:nodule}{Nitrogenase} dimusnahkan oksigen, tetapi bakteroidnya harus bernapas
untuk membuat enam belas ATP tiap N$_{2}$. Korteks bintilnya membentuk sebuah
penghadang pembauran yang membatasi masuknya oksigen, dan \omterm{def:b3:microbiomes:nodule}{leghemoglobin}, pada
kepekatan milimolar, mengikat oksigen yang masuk sehingga kepekatan
bebasnya bertahan dekat \qty{10}{nM} sementara fluks ke bakteroidnya
tinggi. Merah muda adalah oksileghemoglobin; sebuah bintil hijau telah
merombak \omterm{def:b3:microbiomes:nodule}{leghemoglobinnya} --- ia menua dan tidak lagi memfiksasi.
\end{solution}

\begin{solution}{exo:b3:microbiomes:8}
Keduanya gagal: tidak ada bintil dan tidak ada \omterm{def:b3:microbiomes:mycorrhiza}{mikoriza arbuskular}, karena
lonjakan kalsium adalah isyarat bersama \omterm{def:b3:microbiomes:mycorrhiza}{jalur simbiosis bersama}.
Jalur itu karena itu mendahului pembintilan --- ia dibangun bagi
\omterm{def:b3:microbiomes:symbiosis}{simbiosis} jamur empat ratus juta tahun lalu lalu direkrut
kacang-kacangan bagi \omterm{def:b3:microbiomes:nodule}{rizobia} kemudian.
\end{solution}

\begin{solution}{exo:b3:microbiomes:9}
Kiers memaparkan sebagian bintil kedelai kepada sebuah atmosfer argon dan
oksigen, yang di dalamnya \omterm{def:b3:microbiomes:nodule}{rizobianya} tidak dapat memfiksasi nitrogen, sementara
bintil tumbuhan itu selebihnya memfiksasi normal; tumbuhannya memangkas pasokan
oksigen ke bintil yang tidak memfiksasi dan \omterm{def:b3:microbiomes:nodule}{rizobia} bintil itu berkembang biak
sekitar separuhnya saja. Tanpa pembedaan semacam itu, \omterm{def:b3:microbiomes:nodule}{rizobia} yang melewatkan
fiksasi yang mahal itu akan berkembang biak lebih cepat daripada pemfiksasi,
menyebar ke seluruh populasi, dan tumbuhannya akan berakhir dengan merumahkan
bakteri yang tidak memberi apa-apa: \omterm{def:b3:microbiomes:symbiosis}{mutualismenya} akan runtuh.
\end{solution}

\begin{solution}{exo:b3:microbiomes:10}
$\binom{N-N_{i}}{n}/\binom{N}{n} = \prod_{j=0}^{n-1}(N - N_{i} - j)/(N -
j)$, dengan tiap faktornya di bawah $1$; beranjak dari $n$ ke $n + 1$ mengalikan
satu faktor lagi di bawah $1$, sehingga peluang terlewatnya turun dan peluang
hadirnya naik seiring $n$; pada $n = N$ pembilang $\binom{N -
N_{i}}{N}$ nol dan tiap spesies hadir, yang memberi $S$. Bagi $N_{i},
n \ll N$ tiap faktornya kira-kira $1 - N_{i}/N$, hasil kalinya kira-kira $(1 -
N_{i}/N)^{n} \approx e^{-nN_{i}/N} \approx (1 - n/N)^{N_{i}}$ --- yakni
hampiran pencontohan dengan pengembalian.
\end{solution}

\begin{solution}{exo:b3:microbiomes:11}
Seekor jantan adalah jalan buntu bagi simbion yang mengembara lewat telur saja,
sehingga sifat apa pun yang mengubah jantan menjadi betina, membunuhnya demi
saudarinya, atau melumpuhkan betina yang tidak terinfeksi menaikkan
penularan simbionnya. Ketaksesuaian sitoplasma: sperma jantan yang terinfeksi
diubah sehingga zigotnya mati kecuali telurnya mengusung
\emph{Wolbachia} yang sama, yang menyelamatkannya. Betina yang tidak terinfeksi
karena itu kehilangan keturunan yang dikandungnya bersama jantan yang terinfeksi,
sementara betina yang terinfeksi tidak kehilangan apa pun; keuntungannya tumbuh
seiring kekerapan infeksinya, sehingga di atas sebuah ambang infeksinya menyebar
sampai terpancang walaupun betina yang terinfeksi bertelur sedikit lebih sedikit
--- ongkosnya tertimpa bagian anakan betina tak terinfeksi yang musnah.
\end{solution}

\begin{solution}{exo:b3:microbiomes:12}
Endosimbionnya melewati sebuah leher botol beranggota beberapa sel pada tiap
angkatan inang lalu tidak pernah berekombinasi dengan pihak luar: hanyutan
memancangkan mutasi yang agak merugikan, termasuk penghapusan, dan tidak ada
sumber gen pengganti (ratcet Muller). Gen yang produknya
dipasok inangnya, atau yang melayani hidup bebas --- \omterm{def:b3:virology:virus}{selubung}, gerak,
pengaturan, sebagian besar perbaikan --- tidak dirindukan lalu hilang lebih
dulu; mesin penerjemahannya dan gen yang membuat apa yang diperlukan inangnya
(asam amino esensial pada \emph{Buchnera}) bertahan paling lama. Kerabatnya yang
hidup bebas berpopulasi raksasa, berekombinasi, dan kemasukan gen, sehingga
seleksi memelihara genomnya. Pembalikannya memerlukan gen baru dari
luar --- mustahil dalam pengasingan --- sehingga inangnya justru mengganti
simbion yang terkikis dengan yang segar, sebagaimana telah dilakukan beberapa
garis keturunan serangga.
\end{solution}

\begin{solution}{pb:b3:microbiomes:1}
\textbf{1.} $400\times 10^{11} = 4\times 10^{13}$ bakteri, \qty{40}{g}.
\textbf{2.} Sebuah genom bergen \num{4000} kira-kira \qty{4}{Mb}, \qty{4.4}{fg}:
$4\times 10^{13}\times 4.4\times 10^{-15} = \qty{0.18}{g}$ DNA
bakteri, terhadap $3\times 10^{13}\times 6.4\times 10^{-12} = \qty{0.19}{g}$
DNA manusia --- kira-kira sama.
\textbf{3.} $1000\times 0.7\times 4000 + 0.3\times 4000 \approx 2.8\times
10^{6}$ gen berbeda, sekitar $140$ kali lipat milik genom manusia.
\textbf{4.} $30\times 10 = \qty{300}{mmol}$; $\times 0.9 = \qty{270}{kJ}$;
yakni \qty{3}{\%} dari makanannya.
\textbf{5.} Tidak --- \qty{3}{\%} tidak menjelaskan \qty{30}{\%}. Mencit bebas
kuman juga menyerap kurang mangkus, berhormon usus, bersimpanan lemak, dan
bertermogenesis yang berubah, serta berselera beda: \omterm{def:b3:microbiomes:symbiosis}{mikrobiomnya} bekerja
pada fisiologi inangnya, bukan hanya sebagai pemasok bahan bakar.
\textbf{6.} Pembelahan diimbangi pembuangan: sebuah populasi mantap, tiap selnya
tergantikan tiap hari. Sesudah pembunuhan \qty{99.9}{\%}, $4\times 10^{10}$ tersisa
dan sepuluh penggandaan memulihkan cacahnya dalam sekitar sepuluh hari --- tetapi
yang selamat dan yang tercepat tumbuh menguasainya, susunannya
berubah (sebuah \omterm{def:b3:microbiomes:gut}{disbiosis}), dan pada celah itu seorang penyerbu seperti \emph{C.
difficile} dapat mapan.
\textbf{7.} Spesies kecilnya berbagi $0.4$ atas $177$: $p = 0.00226$ masing-masing. $H
= -(0.3\ln 0.3 + 0.2\ln 0.2 + 0.1\ln 0.1 + 0.4\ln 0.00226) = 0.361 +
0.322 + 0.230 + 2.437 = 3.35$; $e^{H} \approx 28$.
\textbf{8.} $D = 0.09 + 0.04 + 0.01 + 177\times (0.00226)^{2} = 0.141$;
$1/D = 7.1$. \omterm{def:b3:microbiomes:diversity}{Indeks Simpson} mengkuadratkan kelimpahannya, sehingga ketiga
spesies yang menguasai itulah yang memutuskannya dan yang jarang hampir tak
terhitung; Shannon memberi bobot lebih kepada spesies yang jarang.
\textbf{9.} $1 - 60/2000 = 0.97$: sekitar $30$ \omterm{def:b3:genomics:ngs}{bacaan} tiap seribu tergolong
spesies yang belum terlihat.
\textbf{10.} Kekayaannya naik seiring ukuran contohnya sementara spesies jarang
terus bermunculan; \num{10000} \omterm{def:b3:genomics:ngs}{bacaan} menjangkau spesies yang dilewatkan
\num{2000} \omterm{def:b3:genomics:ngs}{bacaan}. Jarangkanlah contoh besarnya ke \num{2000} \omterm{def:b3:genomics:ngs}{bacaan} dengan
teoremanya (atau dengan mengambil anak contoh) --- ia mestinya memberi sekitar
$180$ bila masyarakatnya sama --- lalu bandingkanlah pada kedalaman yang sama.
\textbf{11.} $N_{i} = 1$: $1 - \binom{1999}{1000}/\binom{2000}{1000} =
1 - (2000 - 1000)/2000 = 0.5$. $N_{i} = 5$: kira-kira $1 - 0.5^{5} = 0.97$.
\textbf{12.} $e^{2.0} = 7.4$ spesies mujarab; $D \ge 0.5^{2} = 0.25$,
sekitar $0.27$: sebuah masyarakat dengan satu spesies yang menguasai dan
beberapa lusin spesies kecil --- yakni usus yang runtuh dan miskin keanekaan
yang di dalamnya \emph{C.~difficile} berpijak.
\textbf{13.} $150\times 8 = \qty{1200}{kg}$ gula tiap hektare,
yakni \qty{10}{\%} dari \qty{12}{t} yang difotosintesiskan.
\textbf{14.} $150\,000/28 = 5360$ mol N$_{2}$; $16\times 5360 =
\num{85700}$ mol ATP.
\textbf{15.} $\num{85700}/30 = 2860$ mol glukosa, yakni \qty{514}{kg} ---
sekitar \qty{43}{\%} dari gula yang dibayar. Sisanya membangun dan memelihara
bintil dan bakteroidnya, memasok kedelapan elektron tiap N$_{2}$ sebagai
pereduksi, lalu mengasimilasi amonianya.
\textbf{16.} $1200\times 16 = \qty{19200}{MJ}$ bagi \qty{150}{kg} N:
yakni \qty{128}{MJ/kg}, hampir empat kali \qty{35}{MJ/kg} pabriknya --- tetapi
dibayar dengan sinar matahari, di ladang, tanpa bahan bakar.
\textbf{17.} Pemfiksasi $0.9\times 1 = 0.9$; pencurang $0.1\times 0.5 =
0.05$; pencurangnya $0.05/0.95 = 5.3\,\%$ --- terpotong separuh dalam satu musim.
\textbf{18.} Pencurang $0.1\times 1.2^{10} = 0.62$ terhadap pemfiksasi $0.9$:
yakni \qty{41}{\%} sesudah sepuluh musim, yang naik ke arah penguasaan. Sanksi
memutar seleksi melawan pencurangnya; tanpanya \omterm{def:b3:microbiomes:symbiosis}{mutualismenya} terkikis.
\textbf{19.} $10^{6}\times \qty{20}{pg} = \qty{20}{\micro g}$ karbon
tiap cm$^{2}$ tiap hari; \qty{18}{\micro g} ke karangnya.
\textbf{20.} Pernapasannya \qty{15}{\micro g}: lebihnya \qty{3}{\micro g}
tiap cm$^{2}$ tiap hari bagi pertumbuhan, rangka organik pengapuran,
lendir, dan gamet.
\textbf{21.} Dengan \qty{5}{\%} alganya, \qty{0.9}{\micro g} sehari
terhadap \qty{15}{\micro g} yang dinapaskan: sebuah kekurangan \qty{14}{\micro g}
sehari; cadangan \qty{300}{\micro g} bertahan sekitar tiga minggu.
\textbf{22.} Empat minggu pada $+2$ lebih buruk: kerusakan pada fotosistem
alganya naik dengan curam, bukan lurus, terhadap suhunya. Ukurannya
toh berhasil sebagai peringatan karena pemutihan mengumpulkan tekanan yang
menimbun dan kedua tampangnya berbeda kurang daripada ketakpastian pada
ambangnya; ia sebuah indeks empiris, yang ditera pada pemutihan yang teramati.
\textbf{23.} Sebuah karang yang tertekan dapat beralih ke spesies alga yang
lebih tahan panas yang sudah hadir pada kelimpahan rendah, atau memperolehnya dari
airnya sesudah pemutihan; alga yang tahan itu memberi lebih sedikit karbon, sehingga
karangnya tumbuh lebih lambat, dan ketahanan yang diperolehnya satu atau dua
derajat --- kurang daripada pemanasan yang diramalkan --- sementara laju
pemanasannya melampaui penyesuaian diri karangnya sendiri.
\textbf{24.} $10^{10}$ cm$^{2}$ $\times$ \qty{20}{\micro g} $=
\qty{200}{kg}$ karbon sehari, yakni sekitar \qty{73}{t} setahun.
\textbf{25.} Sekitar \qty{3}{\%} makanannya dari peragian
\omterm{def:b3:microbiomes:symbiosis}{mikrobiomnya}; cakupan contoh pengurutannya \qty{97}{\%};
\qty{10}{\%} fotosintesis dibayar bagi nitrogen tanaman kacangnya.
\end{solution}
